% Einheit 2 von 3 – Zielfunktion aus Haupt- und Nebenbedingung (bank/extremalprobleme/e2.jsonl, zusammenbau v0.1)
%% TODO Zeitmarke Einheit 2: Marken-Zeile nicht lesbar
%% TODO Zweigzeile Teil 1: was hier gelernt wird (Satz in Schülersprache aus dem Einheitstitel) – Einheit 2
\einheitenkopf[e2]{Einheit 2 von 3 · Zielfunktion aus Haupt- und Nebenbedingung}
\zweigzeile{FHR}
\setcounter{aufgabe}{10}

%% TODO Titel: Ich-kann-Satz fehlt in der Bank (Platzhalter: Kettenname) – Nr. 11
%% TODO Anweisung: ein Satz über der ersten Teilaufgabe fehlt in der Bank – Nr. 11
\begin{aufgabe}{Zielfunktion aufstellen}
%% TODO \rechenplatz{n}: Schrittzahl des Grundfalls fehlt in der Bank (nur bei fester Schreibform, 2.3 b)
\begin{teile}
\teil Mit 60 m Zaun soll an einer Mauer ein rechteckiges Beet eingezäunt werden; der Zaun bildet drei Seiten, $a$ steht senkrecht zur Mauer, $b$ parallel. Die Fläche soll möglichst groß werden. Welche Gleichung ist die Nebenbedingung? \kreuz{$A = a \cdot b$} \\ \kreuz{$2a + b = 60$} \\ \kreuz{$2a + 2b = 60$}
\teil Ein rechteckiges Gehege wird mit 40 m Zaun vollständig umzäunt; seine Fläche soll möglichst groß werden. Notiere Haupt- und Nebenbedingung und stelle die Nebenbedingung nach $b$ um.
\teil An einer Hauswand liegt ein rechteckiger Hundeauslauf; für die drei übrigen Seiten stehen 30 m Begrenzung zur Verfügung, die vollständig verwendet wird. $a$ ist die Seite senkrecht zur Wand, $b$ die Seite parallel zur Wand. Berechne den Flächeninhalt für $a = 6$ m und für $a = 10$ m. \hfill (FHR 2025)
\teil Für ein rechteckiger Spielplatz an einem Bach stehen für die drei Seiten ohne Bach, Mauer oder Stall 70 m Zaun zur Verfügung, die vollständig verwendet werden; $a$ ist die Seite senkrecht dazu, $b$ die Seite parallel. Gib eine Funktion an, die den Flächeninhalt in Abhängigkeit von $a$ beschreibt, mit Definitionsbereich. \hfill (FHR 2025) \\ A(a) = \leerfeld
\teil Die Graphen von $f(x) = 5 - 0{,}2x^2$ und $-f$ begrenzen ein Gehäuse. Darin liegt ein Rechteck mit den Ecken $(-x | -f(x))$, $(x | -f(x))$, $(x | f(x))$ und $(-x | f(x))$, $0 < x < 5$. Weise nach, dass sein Flächeninhalt $A(x) = 20x - 0{,}8x^3$ beträgt.
\teil Der Ursprung und der Punkt $P(u | f(u))$ mit $u > 0$ auf dem Graphen von $f(x) = (x + 3) \cdot e^{-0{,}5x}$ sind gegenüberliegende Ecken eines achsenparallelen Rechtecks. Stelle die Zielfunktion für den Flächeninhalt in Abhängigkeit von $u$ auf. \hfill (Abitur 2022 GK) \\ A(u) = \leerfeld
\teil Für $k > 0$ ist $f_k(x) = 2x \cdot e^{-kx}$. Für $a > 0$ bilden die Punkte $(0 | 0)$, $(a | 0)$ und $(a | f_k(a))$ ein Dreieck. Begründe, dass sein Flächeninhalt $a^2 \cdot e^{-ka}$ beträgt. \hfill (Abitur 2020 LK)
\end{teile}
\end{aufgabe}

%% TODO Titel: Ich-kann-Satz fehlt in der Bank (Platzhalter: Kettenname) – Nr. 12
%% TODO Anweisung: ein Satz über der ersten Teilaufgabe fehlt in der Bank – Nr. 12
\begin{aufgabe}{Zielfunktion aufstellen – Fehler finden, Begründen, Anwendung}
\begin{teile}
\teil Finde den Fehler. An einer Wand stehen für drei Seiten eines Rechtecks 50 m Material zur Verfügung; $a$ steht senkrecht zur Wand, $b$ parallel. Gesucht ist der Flächeninhalt als Funktion von $a$. Rechnung: $2a + b = 50$, also $a = 25 - 0{,}5b$ und $A(b) = (25 - 0{,}5b) \cdot b$.
\teil Begründe, warum die Zielfunktion nur noch eine Variable enthalten darf.
\teil Eine zylinderförmige Konservendose soll 1 Liter fassen ($1\,000$ cm³). Stelle ihre Oberfläche als Funktion des Radius $r$ (in cm) auf und berechne, wie viel Blech eine Dose mit $r = 5$ cm braucht.
\end{teile}
\end{aufgabe}
